\begin{figure}[htbp]
\centering
\resizebox{\textwidth}{!}{%
\begin{tikzpicture}[
  font=\sffamily\small,
  block/.style={draw=VIUDark!70,fill=VIULight,rounded corners=4pt,
    align=center,text width=3.0cm,minimum height=1.0cm,inner sep=5pt},
  core/.style={draw=VIUOrange!90!black,fill=VIUOrange!8,rounded corners=4pt,
    align=center,text width=3.0cm,minimum height=1.0cm,inner sep=5pt,
    font=\sffamily\small\bfseries},
  flow/.style={-{Latex[length=2.3mm]},line width=0.9pt,draw=VIUOrange!90!black},
  feed/.style={-{Latex[length=2.0mm]},line width=0.8pt,draw=VIUDark!60}
]
\node[core]  (load) at (0,1.2) {Carga r\'igida\\estado y energ\'ia};
\node[block] (cost) at (3.8,1.2) {Coste\\anisotr\'opico};
\node[core]  (solver) at (7.6,1.2) {Flujo continuo\\HJB--FPK};
\node[block] (field) at (11.4,1.2) {Campo ideal\\direcci\'on local};

\node[block] (rho) at (0,-1.05) {Densidad\\fase colectiva};
\node[block] (port) at (3.8,-1.05) {Puerto mec\'anico\\fuerza y torque};
\node[core]  (safe) at (7.6,-1.05) {Filtro seguro\\MF-CBF / QP};
\node[block] (uni) at (11.4,-1.05) {Uniciclo\\mando motor};

\draw[flow] (load) -- (cost);
\draw[flow] (cost) -- (solver);
\draw[flow] (solver) -- (field);
\draw[flow] (field) -- (uni);
\draw[flow] (uni) -- (safe);
\draw[flow] (safe) -- (port);
\draw[flow] (port) -- (load);
\draw[feed] (rho) -- (solver);
\draw[feed] (load) -- (rho);
\draw[feed] (rho) -- (port);

\node[draw=VIUGray!55,fill=VIUWhite,rounded corners=4pt,
  align=center,text width=11.0cm,inner sep=5pt,font=\sffamily\scriptsize,
  below=0.6cm of port] {Lectura: el flujo continuo propone una direcci\'on; cada AMR la discretiza, la filtra y devuelve esfuerzo a la carga.};
\end{tikzpicture}}
\caption{Arquitectura matemática del transporte cooperativo amorfo. La carga modifica el
coste de campo medio; la densidad resuelve un flujo continuo; cada AMR discretiza ese
flujo mediante un punto adelantado y un filtro CBF; el contacto resultante vuelve a
inyectar fuerza en el puerto port-Hamiltoniano de la carga.}
\viuownsource
\label{fig:transporte-amorfo-stack}
\end{figure}
